\section{Ontological Knowledge Graph in Neo4j}
\label{sec:knowledge_graph}

The deterministic reasoning of Omni-IPS depends on a strongly typed, graph-theoretic representation of domain axioms and facts. We utilize \textbf{Neo4j} as our logical Graph Database, which acts as the structured grounding layer (ontology hook) of the system. By structuring theorems and assertions as a graph, we can represent dependencies, input-output logic flows, and taxonomic classifications directly in the database schema~\cite{pan2024unifying}.

\subsection{Ontology Schema Design}
The knowledge graph is modeled using two primary node labels and two directed relationship types, forming an bipartite graph structure of rules and facts:
\begin{itemize}
    \item \textbf{Nodes}:
        \begin{itemize}
            \item \texttt{Fact}: Represents a concrete physical or mathematical assertion. It possesses the properties: \texttt{id} (unique key), \texttt{value} (exact string expression, e.g., ``BaSO4''), \texttt{domain} (chemistry, geometry, or algebra), and optional metadata attributes (e.g., solubility class, state of aggregation).
            \item \texttt{Rule}: Represents an abstract logical transition or theorem. It contains: \texttt{id} (e.g., ``geo\_pythagoras''), \texttt{name} (human-readable theorem name), \texttt{domain}, and \texttt{description} explaining the rule's mechanical transformation.
        \end{itemize}
    \item \textbf{Relationships}:
        \begin{itemize}
            \item \texttt{(f:Fact)-[:HAS\_INPUT]->(r:Rule)}: Declares that fact $f$ is a prerequisite (antecedent) required to fire rule $r$.
            \item \texttt{(r:Rule)-[:HAS\_OUTPUT]->(f:Fact)}: Declares that fact $f$ is a logical conclusion (consequent) generated when rule $r$ fires.
        \end{itemize}
\end{itemize}

The resulting schema is visualized in Figure~\ref{fig:neo4j_kb} (derived from the actual database visualization in \texttt{assets/neo4j\_knowledge\_base.png}), which displays the bipartite mapping of mathematical theorems and elements.

\begin{figure}[h]
    \centering
    \includegraphics[width=0.85\textwidth]{neo4j_knowledge_base.png}
    \caption{Bipartite Neo4j Knowledge Graph visualization displaying the connection between Rule nodes (theorems) and Fact nodes via \texttt{HAS\_INPUT} and \texttt{HAS\_OUTPUT} relationships.}
    \label{fig:neo4j_kb}
\end{figure}

\subsection{Commutative Relation Canonicalization}
A major issue in propositional reasoning is argument permutation in symmetric relations. For example, in geometry, the parallel lines assertion:
\begin{equation*}
    \texttt{Parallel}(AB, CD) \equiv \texttt{Parallel}(CD, AB)
\end{equation*}
Similarly, segment congruence:
\begin{equation*}
    \texttt{Congruent}(AB, CD) \equiv \texttt{Congruent}(CD, AB) \equiv \texttt{Congruent}(BA, CD)
\end{equation*}
Because our core reasoning engine uses exact string matching (propositional equality), treating these permutations as separate strings would lead to reasoning failure unless every permutation was explicitly added to working memory. 

To prevent this exponential state explosion, we enforce **Commutative Canonicalization** inside the domain-specific parser during Neo4j ingestion. The set of commutative relations $\mathcal{C}$ is defined as:
\begin{equation}
    \mathcal{C} = \{ \texttt{Congruent}, \texttt{Similar}, \texttt{Parallel}, \texttt{Intersect}, \texttt{Equal} \}
\end{equation}
When a fact $R(A_1, A_2, \dots, A_n)$ is parsed, if $R \in \mathcal{C}$, the arguments are lexicographically sorted:
\begin{equation}
    \text{canonical\_args} = \text{sort}(A_1, A_2, \dots, A_n)
\end{equation}
For example, a parsed fact \texttt{Parallel(CD, AB)} automatically maps to the canonical value \texttt{Parallel(AB, CD)} in both the Neo4j database and the working memory, ensuring robust unification.

\subsection{Cypher Rule Loading}
During system query execution, the FastAPI gateway queries all rules associated with the active domain from Neo4j. To guarantee robustness, we query both structural database relationships and fallback array properties using the following optimized Cypher query:

\begin{lstlisting}[language=SQL, style=jsonstyle, caption={Cypher query executed by the backend to retrieve domain rules and their connected input/output facts.}]
MATCH (r:Rule) WHERE r.domain = $domain
OPTIONAL MATCH (f_in:Fact)-[:HAS_INPUT]->(r)
OPTIONAL MATCH (r)-[:HAS_OUTPUT]->(f_out:Fact)
WITH r, 
     CASE WHEN r.inputs IS NOT NULL 
          THEN r.inputs 
          ELSE collect(DISTINCT f_in.value) 
     END AS inputs,
     CASE WHEN r.outputs IS NOT NULL 
          THEN r.outputs 
          ELSE collect(DISTINCT f_out.value) 
     END AS outputs
RETURN r.id AS id, 
       r.name AS name, 
       inputs, 
       outputs, 
       r.description AS description;
\end{lstlisting}
This query retrieves rules dynamically, assembling antecedents and consequents into neat lists of string values matching our Pydantic rule schema.
